\scanpage{60}
% archive sheet number: 59

% L1–3, proposition
\uline{Proposition} $F$ faisceau abélien sur un espace top. (ordinaire) $X$. Conditions équivalentes

% L4, condition i
(i) Pour tt ouvert $U$ de $X$, $H^1(U,F)=0$

% L5, condition i bis
(\deleted{\unsure{ii}}\unsure{i bis|ii}) Pour tt ouvert $U$ de $X$, $H^i(U,F)=0$ pour $i>0$.

% L6, condition ii
(\deleted{\unsure{iii}}\unsure{ii|iii}) \inserted{[pour $X$ noeth]} Pour deux ouverts $U,V$ de $F$,

% L7, two-open formula
\[
H^1((U,V);F)=H^0(U\cap V,F)/\operatorname{Im}\big(H^0(U,F)\times H^0(V,F)\big)
\]

% L8, continuation
est nul.

% L9–10, condition ii bis
(ii bis) Pour tt ouvert $U$ et tt recouvrement \inserted{ouvert} $\underline{U}=(U_i)$ de $U$, on a

% L11, cover cohomology
\[
H^i(\underline{U},F)=0\qquad\text{pour}\quad i>0.
\]

% L12–14, first equivalence
(i bis) $\Longrightarrow$ (ii bis) par suite spectrale de Leray, est bien connu (petit argument de Cartan en une direction)

% L15–16, two implications
\[
\begin{aligned}
(\mathrm{i\ bis})&\Longrightarrow(\mathrm{i})\\[-2pt]
(\mathrm{ii\ bis})&\Longrightarrow(\mathrm{ii})
\end{aligned}
\quad\text{trivial}
\]

% L17–21, final argument
(i) $\Longleftrightarrow$ (ii) En effet $H^1(\deleted{\unsure{\{} }U\cap V,F)\subset H^1(U\cup V,F)$, d'où (i) $\Rightarrow$ (ii), pour (ii) $\Longrightarrow$ (i) on \deleted{\unsure{prend}} considère $\unsure{\zeta|\xi}\in H^1(U,F)$, \deleted{\unsure{on note}} considère un ouvert $U'\subset U$ \uline{maximal} \inserted{(N.B. $X$ noeth)} où $\zeta$ soit nul, je dis que $U'=U$, \unsure{on regarde un|en regardant}.
